\section[A basis for $\Delta $]{A basis for $\boldsymbol{\Delta}$}

 In this section we display a basis for the $\mathbb F$-vector
space $\Delta $.

\begin{Theorem}
\label{thm:basis}
The following is a basis for the $\mathbb F$-vector
space $ \Delta $.
\begin{gather}
\label{eq:basis}
A^iB^jC^k \alpha^r\beta^s\gamma^t, \qquad   i,j,k,r,s,t \geq 0.
\end{gather}
\end{Theorem}

\begin{proof}
We invoke  Bergman's Diamond Lemma
\cite[Theorem~1.2]{berg}.
Consider the symbols
\begin{gather}
\label{eq:letters}
A, \quad B, \quad C, \quad \alpha,\quad \beta, \quad \gamma.
\end{gather}
For an integer $n \geq 0$, by a {\it $\Delta $-word of length $n$}
we mean a sequence
 $x_1x_2\cdots x_n$ such that~$x_i$ is listed in
(\ref{eq:letters}) for $1 \leq i \leq n$.
We interpret the
$\Delta $-word
of length zero to be the multiplicative
identity in $\Delta $.
Consider a
 $\Delta $-word
$w=x_1x_2\cdots x_n$.
By an {\it inversion} for $w$
we mean an ordered pair of integers $(i,j)$
such that
$1\leq i<j\leq n$ and
$x_i$ is strictly to the right of
$x_j$ in the list~(\ref{eq:letters}).
For example $CABA$ has 4 inversions
and
$CB^2A$ has 5 inversions.
A $\Delta $-word is called {\it reducible} whenever it has at least one
inversion, and {\it irreducible} otherwise.
The list
(\ref{eq:basis}) consists of the irreducible $\Delta $-words.
For each integer $n\geq 0$ let
$W_n$ denote the set of $\Delta $-words that have
length~$n$.
Let $W=\cup_{n=0}^\infty W_n$ denote the
set of all
$\Delta $-words.
We now def\/ine a partial
order~$<$ on~$W$.
The def\/inition has two aspects.
$(i)$~For all integers  $n> m\geq 0$,
every word in~$W_m$ is less than
every word in~$W_n$, with respect to $<$.
$(ii)$~For an integer $n\geq 0$
the restriction of $<$ to $W_n$ is
described as follows.
Pick $w, w' \in
 W_n$ and write
 $w=x_1x_2\cdots x_n$.
We say that $w$ {\it covers} $w'$ whenever
there exists an integer~$j$ $(2 \leq j \leq n)$
such that $(j-1,j)$ is an inversion for~$w$,
and
$w'$ is obtained from $w$ by interchanging
$x_{j-1},x_j$.
In this case
$w'$ has one fewer inversions than $w$.
Therefore the transitive closure of the covering
relation on $W_n$ is a partial order on $W_n$, and this is
the restriction of
 $<$ to $W_n$.
We have now def\/ined a partial order $<$ on $W$.
By construction this partial order is
a semi-group partial order
\cite[p.~181]{berg} and satisf\/ies the
descending chain condition
\cite[p.~179]{berg}.
By Def\/inition~\ref{def:uaw} and
Def\/inition~\ref{def:abc}
the def\/ining relations for $\Delta $ can be expressed
as follows:
\begin{gather*}
 BA  =  q^2AB+q\big(q^2-q^{-2}\big)C-q\big(q-q^{-1}\big)\gamma,
\\
 CB  =  q^2BC+q\big(q^2-q^{-2}\big)A-q\big(q-q^{-1}\big)\alpha,
\\
 CA  =  q^{-2}AC+q^{-1}\big(q^{-2}-q^2\big)B-q^{-1}\big(q^{-1}-q\big)\beta,
\\
 \alpha A = A \alpha,
  \qquad
\alpha B = B \alpha,
  \qquad
\alpha C = C \alpha,
\\
 \beta A = A \beta,
  \qquad
\beta B = B \beta,
  \qquad
\beta C = C \beta,
\\
 \gamma A = A \gamma,
  \qquad
\gamma B = B \gamma,
  \qquad
\gamma C = C \gamma,
\\
 \beta \alpha = \alpha \beta,
  \qquad
\gamma \beta = \beta \gamma,
  \qquad
\gamma \alpha = \alpha \gamma.
\end{gather*}
The above equations give reduction rules
for $\Delta$-words,
as we now explain.
Let $w=x_1x_2\cdots x_n$ denote a  reducible $\Delta $-word.
Then there exists an integer $j$ $(2 \leq j \leq n)$
such that $(j-1,j)$ is an inversion for~$w$.
In the above list of equations,
there exists an equation
with $x_{j-1}x_j$ on the left-hand side;
in~$w$ we eliminate $x_{j-1}x_j$ using this equation
and thereby express $w$ as a linear combination of
$\Delta $-words, each less than $w$ with respect to $<$.
Therefore
the reduction rules are compatible
with $<$ in the sense of Bergman \cite[p.~181]{berg}.
In order to employ the Diamond Lemma,
we must show that
the ambiguities
are resolvable in the sense of Bergman
\cite[p.~181]{berg}.
There are potentially two kinds of ambiguities; inclusion
ambiguities and overlap ambiguities
\cite[p.~181]{berg}.
For the present example there are no inclusion ambiguities.
The only nontrivial overlap ambiguity
involves the word
$CBA$. This word can be reduced
in two ways; we could evaluate $CB$ f\/irst or we could
evaluate $BA$ f\/irst.
 Either way, after a three-step reduction
we obtain the same result, which is
\begin{gather*}
 q^{-1}CBA = qABC
+\big(q^2-q^{-2}\big)A^2
-\big(q^2-q^{-2}\big)B^2
+\big(q^2-q^{-2}\big)C^2
\\
\phantom{q^{-1}CBA =}{}
- \big(q-q^{-1}\big)A \alpha
+ \big(q-q^{-1}\big)B \beta
- \big(q-q^{-1}\big)C \gamma .
\end{gather*}
Therefore the
overlap ambiguity $CBA$ is re\-sol\-vable.
We conclude that every  ambiguity is
re\-sol\-vable, so
by the Diamond Lemma
\cite[Theorem~1.2]{berg}
the elements
(\ref{eq:basis}) form a basis for
$\Delta $.
\end{proof}

 On occasion we
wish to discuss the coef\/f\/icients
when an element of $\Delta$ is written as a linear
combination of the elements
(\ref{eq:basis}). To facilitate this discussion
we def\/ine a bilinear form
$(\,,) : \Delta \times \Delta \to \mathbb F$
such that $( u,v) = \delta_{u,v}$
for all elements $u$, $v$ in the basis
(\ref{eq:basis}). In other words
the basis
(\ref{eq:basis}) is
orthonormal with respect to
$( \,,) $.
 Observe that
$(\,,)$ is symmetric.
For $u \in \Delta$,
\begin{gather}
u = \sum \big(u, A^iB^jC^k\alpha^r\beta^s\gamma^t \big)
 A^iB^jC^k\alpha^r\beta^s\gamma^t,
\label{eq:bilsum}
\end{gather}
where the sum is over all elements
 $A^iB^jC^k\alpha^r\beta^s\gamma^t$ in the basis
(\ref{eq:basis}).

\begin{Definition}
\label{def:contrib}
\rm
Let $u \in \Delta $. A given element
 $A^iB^jC^k\alpha^r\beta^s\gamma^t$ in the basis
(\ref{eq:basis}) is said to
{\it contribute to
  $u$} whenever $(u,
 A^iB^jC^k\alpha^r\beta^s\gamma^t )
  \not=0$.
\end{Definition}







\section[A filtration of $\Delta$]{A f\/iltration of $\boldsymbol{\Delta}$}


In this section we obtain
a f\/iltration of $\Delta $
which
is related to the basis from Theorem~\ref{thm:basis}.
 This f\/iltration will be useful
when we investigate the center  $Z(\Delta)$ later in the paper.


We recall some notation.
For subspaces $H$, $K$ of $\Delta $ def\/ine
$HK={\rm Span}\lbrace hk | h \in H, \; k \in K\rbrace$.

\begin{Definition}
\label{def:ui}
We def\/ine subspaces
$\lbrace \Delta_n\rbrace_{n=0}^\infty$
of $\Delta $ such that
\begin{gather*}
\Delta_0=\mathbb F 1,
  \qquad
\Delta_1 =\Delta_0+{\rm Span}\lbrace A,B,C,\alpha,\beta,\gamma\rbrace,
 \qquad
\Delta_n=\Delta_1 \Delta_{n-1}, \qquad n=1,2,\ldots
\end{gather*}
\end{Definition}

\begin{Lemma}
\label{lem:filt}
The following $(i)$--$(iii)$ hold.
\begin{enumerate}\itemsep=0pt
\item[$(i)$]
 $\Delta_{n-1}\subseteq \Delta_n$ for  $n\geq 1$.
\item[$(ii)$]
$\Delta=\cup_{n=0}^\infty \Delta_n$.
\item[$(iii)$]
$\Delta_m \Delta_n = \Delta_{m+n} $ for
 $m,n\geq 0$.
\end{enumerate}
\end{Lemma}

\begin{proof}
$(i)$ Since $\Delta_n=\Delta_1 \Delta_{n-1}$
and $1 \in \Delta_1$.
$(ii)$ Since $A$, $B$, $C$, $\alpha$, $\beta$, $\gamma $ generate~$\Delta$.
$(iii)$ Each side is equal to
$(\Delta_1)^{m+n}$.
\end{proof}

By Lemma~\ref{lem:filt} the sequence
$\lbrace \Delta_n\rbrace_{n=0}^\infty$ is a f\/iltration of~$\Delta$ in the sense of
\cite[p.~202]{carter}.


\begin{Lemma}
\label{lem:almostcom}
Each of the following is contained in $\Delta_1$:
\begin{gather*}
qAB-q^{-1}BA,
\qquad
qBC-q^{-1}CB,
\qquad
qCA-q^{-1}AC.
\end{gather*}
\end{Lemma}

\begin{proof}
Each of the three expressions is a linear combination of
$A$, $B$, $C$, $\alpha$, $\beta$, $\gamma$ and these are contained in~$\Delta_1$.
\end{proof}

\begin{Theorem}\label{lem:unbasis}
For all integers $n\geq 0$ the following is
a basis for the~$\mathbb F$-vector space~$\Delta_n$:
\begin{gather}
\label{eq:basislist}
A^iB^jC^k\alpha^r \beta^s \gamma^t, \qquad
i,j,k,r,s,t\geq 0, \qquad
i+j+k+r+s+t \leq n.
\end{gather}
\end{Theorem}

\begin{proof}
The elements~(\ref{eq:basislist}) are linearly independent by
Theorem
\ref{thm:basis}. We show that the
elements~(\ref{eq:basislist}) span $\Delta_n$.
We will use induction on $n$. Assume $n\geq 2$; otherwise
the result holds by Def\/inition~\ref{def:ui}.
By Def\/inition~\ref{def:ui} $\Delta_n$ is
spanned by the set of elements of the form
$x_1x_2\cdots x_n$
where $x_i \in \lbrace 1, A,B,C,\alpha,\beta,\gamma\rbrace$
for $1 \leq i \leq n$.
Therefore $\Delta_n$ is spanned by the set of elements
of the form
$x_1x_2\cdots x_m$ where $0 \leq m\leq n$ and
$x_i \in \lbrace A,B,C,\alpha,\beta,\gamma\rbrace$
for $1 \leq i \leq m$.
Therefore $\Delta_n$ is spanned by $\Delta_{n-1}$ together with
the set of elements
of the form
$x_1x_2\cdots x_n$ where
$x_i \in \lbrace A,B,C,\alpha,\beta,\gamma\rbrace$
for $1 \leq i \leq n$.
Consider such an element
$x_1x_2\cdots x_n$.
By Lemma~\ref{lem:almostcom} and since
each of $\alpha$, $\beta$, $\gamma$ is central,
we f\/ind that
for $2 \leq j \leq n$,
\begin{gather*}
x_1\cdots  x_{j-1}x_j\cdots  x_n \in
\mathbb F x_1\cdots x_{j}x_{j-1}\cdots  x_n + \Delta_{n-1}.
\end{gather*}
By the above comments
$\Delta_n$ is spanned by $\Delta_{n-1}$ together with the
set
\begin{gather*}
A^iB^jC^k\alpha^r \beta^s \gamma^t, \qquad
i,j,k,r,s,t\geq 0, \qquad
i+j+k+r+s+t = n.
\end{gather*}
By this and induction
$\Delta_n$ is spanned by the elements~(\ref{eq:basislist}).
We have shown that
the elements~(\ref{eq:basislist}) form a basis  for the $\mathbb F$-vector space~$\Delta_n$.
\end{proof}

  Let $V$ denote a vector space over
$\mathbb F$ and let $U$ denote a subspace of~$V$.
By a {\it complement of~$U$ in~$V$} we mean
a subspace $U'$ of $V$ such that
$V=U+U'$ (direct sum).


\begin{Corollary}
\label{lem:com}
For all integers $n\geq 1$ the following is a basis
for a complement of $\Delta_{n-1}$ in $\Delta_n$:
\begin{gather*}
%\label{eq:combas}
A^iB^jC^k\alpha^r \beta^s \gamma^t, \qquad
i,j,k,r,s,t\geq 0, \qquad
i+j+k+r+s+t = n.
\end{gather*}
\end{Corollary}
\begin{proof}
Use Theorem
\ref{lem:unbasis}.
\end{proof}



\section[The Casimir element $\Omega$]{The Casimir element $\boldsymbol{\Omega}$}

We turn our attention to the center $Z(\Delta)$.
In this section we discuss a certain
 element $\Omega \in Z(\Delta)$
 called the Casimir element.
The name is
motivated by
\cite[equation~(1.3)]{Zhidd}.
 In Section 7
we will use~$\Omega$ to describe
$Z(\Delta)$.
We acknowledge that the results of this section are
extensions of
\cite[Lemma~1]{posta2},
\cite[Section~6]{koro},
\cite[equation~(1.3)]{Zhidd}.

\begin{Lemma}
\label{lem:six}
The following elements of $\Delta $ coincide:
\begin{gather*}
qABC + q^2A^2+q^{-2}B^2+q^2C^2-q A\alpha -q^{-1} B\beta -q C\gamma,
\\
qBCA +
q^2A^2 +
q^2B^2+q^{-2}C^2
-q A\alpha
-qB\beta -q^{-1}C\gamma,
\\
qCAB +
q^{-2}A^2+q^2B^2
+q^2 C^2
-q^{-1} A\alpha -q B \beta
-q C\gamma,
\\
q^{-1}CBA + q^{-2}A^2+q^{2}B^2+q^{-2}C^2-q^{-1}A\alpha -q B\beta -
q^{-1} C\gamma,
\\
q^{-1}ACB +
q^{-2}A^2 +
q^{-2}B^2+q^{2}C^2
-q^{-1} A\alpha
-q^{-1}B\beta -q C\gamma,
\\
q^{-1}BAC +
q^{2}A^2+q^{-2}B^2
+q^{-2}C^2
-q A\alpha
-q^{-1} B \beta
-q^{-1} C\gamma.
\end{gather*}
We denote this common element by $\Omega$.
\end{Lemma}

\begin{proof}
Denote the displayed sequence of elements by
$\Omega^+_B$,
$\Omega^+_C$,
$\Omega^+_A$,
$\Omega^-_B$,
$\Omega^-_C$,
$\Omega^-_A$.
The automorphism
$\rho$ cyclically permutes
$\Omega^+_A$,
$\Omega^+_B$,
$\Omega^+_C$
and cyclically permutes
$\Omega^-_A$, $\Omega^-_B$, $\Omega^-_C$.
The element
$\Omega^+_B-\Omega^-_C$
is
equal to $(q-q^{-1})A$ times
\begin{gather}
\label{eq:need0}
\big(q+q^{-1}\big)A + \frac{qBC-q^{-1}CB}{q-q^{-1}} -\alpha.
\end{gather}
The element
(\ref{eq:need0}) is zero by Def\/inition
\ref{def:abc} so
$\Omega^+_B=\Omega^-_C$.
In this equation we apply $\rho$ twice to get
$\Omega^+_C=\Omega^-_A$
and
$\Omega^+_A=\Omega^-_B$.
The element
$\Omega^+_B-\Omega^-_A$ is equal to
\begin{gather}
\label{eq:need02}
\big(q+q^{-1}\big)C + \frac{qAB-q^{-1}BA}{q-q^{-1}} -\gamma
\end{gather}
times $(q-q^{-1})C$.
The element
(\ref{eq:need02}) is zero by Def\/inition
\ref{def:abc} so
$\Omega^+_B=\Omega^-_A$.
Applying $\rho$ twice we get
$\Omega^+_C=\Omega^-_B$
and
$\Omega^+_A=\Omega^-_C$.
By these comments
$\Omega^+_B$,
$\Omega^+_C$,
$\Omega^+_A$,
$\Omega^-_B$,
$\Omega^-_C$,
$\Omega^-_A$ coincide.
\end{proof}

\begin{Theorem}
\label{lem:center}
The element $\Omega$ from
Lemma~{\rm \ref{lem:six}} is central in $\Delta $.
\end{Theorem}

\begin{proof}
We f\/irst show $\Omega A = A \Omega $.
We will work with the equations~(\ref{eq:u2}),
(\ref{eq:u3}) from Def\/inition~\ref{def:abc}. Consider the equation which is
$qC $ times~(\ref{eq:u2}) plus~(\ref{eq:u2}) times $q^{-1}C$
minus $\gamma$ times~(\ref{eq:u2})
plus $\beta$ times~(\ref{eq:u3})
minus
$q^{-1}B$ times~(\ref{eq:u3})
minus~(\ref{eq:u3}) times~$qB$.
After some cancellation this equation yields
$\Omega^+_B A- A\Omega^+_C =0$, where
$\Omega^+_B$, $\Omega^+_C$
are from the proof of Lemma
\ref{lem:six}.
Therefore
$\Omega A= A\Omega $.
 One similarly
f\/inds
$
\Omega B = B \Omega$ and
$
\Omega C = C \Omega$.
The elements $A$, $B$, $C$ generate~$\Delta $ so~$\Omega$ is central in~$\Delta$.
\end{proof}
